\chapter{Tauberian Theorems}

\begin{center}\itshape
A series can have a sum without converging. Abel, Ces\`aro and Borel all found ways to assign one, and each way is faithful to the ordinary sum when the
ordinary sum exists. The hard question is the converse: when does the assigned sum force the series to converge? The answers are the Tauberian theorems,
and the best of them turn a statement about a transform into a statement about the terms. One of them proves the prime number theorem.
\end{center}

In 1826, in the first volume of Crelle's new journal, a twenty-four-year-old Norwegian proved that a power series is continuous up to a point of the
boundary where it still converges. Niels Henrik Abel was not interested in continuity for its own sake. He was interested in the fact that his
contemporaries added divergent series as if they were numbers, and he wanted to know when that was allowed; in a footnote of the same paper he politely
demolished a theorem of Cauchy on the continuity of sums of series, with a counterexample that is now in every course. Abel is often quoted as calling
divergent series an invention of the devil; the sentence is not easy to trace in his published writings, but his letters of 1826 complain, in plain words,
that in higher analysis very few propositions are proved with convincing rigour.

Abel's theorem says: if $\sum a_n$ converges to $s$, then $f(x)=\sum a_nx^n\to s$ as $x\to1^-$. The converse is false, since $\sum(-1)^n$ does not converge
although $f(x)=\frac1{1+x}\to\frac12$. So a limit of $f$ at $1$ is a weaker piece of information than convergence, and the game is to find the extra
hypothesis that closes the gap. In 1897 Alfred Tauber found the first one: $na_n\to0$. It is a modest theorem with an immodest future, because Hardy and
Littlewood saw that the pattern --- an asymptotic statement about a transform, plus a one-sided or smallness condition on the terms, yields an asymptotic
statement about the terms --- is everywhere in analysis, and they named the pattern after Tauber.

This chapter proves four Tauberian theorems of increasing power. Tauber's own, in one page. Karamata's, which reads the asymptotics of a measure off the
asymptotics of its Laplace transform, and whose only tool is the Weierstrass approximation theorem. Wiener's, the deepest of the four, which says that one
function of $L^1(\R)$ whose Fourier transform never vanishes generates the whole of $L^1(\R)$ by translations: Chapter~16 promised it, and it is the reason
the Fourier transform of a nowhere-zero function matters. And Wiener--Ikehara, which converts a Dirichlet series with a single pole into the asymptotics of
its coefficients, and gives the prime number theorem with no zero-free region, no contour shifting and no explicit formula: only the non-vanishing of
$\zeta$ on the line $\operatorname{Re}s=1$. We end by counting the sums of two squares, in a logarithmic form that Karamata's theorem reaches in half a page
and whose unweighted version needed Landau's full machinery.

\section{Abel and Ces\`aro: regular ways to sum}

All the summation methods we need are weighted averages of the partial sums, so we prove the one lemma they share. Given $a_0,a_1,\dots$, write
$s_n=a_0+\dots+a_n$ for the partial sums and $\sigma_n=\frac{s_0+\dots+s_n}{n+1}$ for their averages.

\begin{lemma}[regular weights]\label{lem:regular}
For $x\in\left]0,1\right[$ and $n\ge0$ let $w_n(x)\ge0$ with $\sum_{n\ge0}w_n(x)=1$, and suppose $w_n(x)\to0$ as $x\to1^-$ for each fixed $n$. If
$s_n\to s$, then $\sum_{n\ge0}w_n(x)s_n\to s$ as $x\to1^-$.
\end{lemma}
\begin{proof}
The sequence is bounded, say $|s_n|\le M$, so the series converges absolutely. Given $\varepsilon>0$, choose $N$ with $|s_n-s|<\varepsilon$ for $n\ge N$. Since the
weights sum to $1$,
\[
\Bigl|\sum_nw_n(x)s_n-s\Bigr|=\Bigl|\sum_nw_n(x)(s_n-s)\Bigr|\le(M+|s|)\sum_{n<N}w_n(x)+\varepsilon.
\]
The first term is a fixed finite sum of quantities tending to $0$, so it is below $\varepsilon$ for $x$ close to $1$.
\end{proof}

\begin{theorem}[Abel, 1826]\label{thm:abel}
If $\sum_{n\ge0}a_n$ converges to $s$, then $f(x)=\sum_{n\ge0}a_nx^n$ converges for $|x|<1$ and $f(x)\to s$ as $x\to1^-$.
\end{theorem}
\begin{proof}
The $a_n$ are bounded, so the series converges for $|x|<1$. For $0<x<1$, Abel summation (or multiplying out) gives
\[
(1-x)\sum_{n\ge0}s_nx^n=\sum_{n\ge0}s_nx^n-\sum_{n\ge1}s_{n-1}x^n=\sum_{n\ge0}a_nx^n=f(x),
\]
all series being absolutely convergent. So $f(x)=\sum_nw_n(x)s_n$ with $w_n(x)=(1-x)x^n\ge0$, which sums to $1$ and tends to $0$ for each fixed $n$. Apply
Lemma~\ref{lem:regular}.
\end{proof}

\begin{corollary}\label{cor:cesaroabel}
If $\sigma_n\to\sigma$, then $f(x)\to\sigma$ as $x\to1^-$. A Ces\`aro-summable series is Abel-summable, with the same sum.
\end{corollary}
\begin{proof}
Write $S(x)=\sum_ns_nx^n$, absolutely convergent for $|x|<1$ because $s_n=O(n)$. Multiplying by $\frac1{1-x}$ sums the coefficients, so
$\sum_n(n+1)\sigma_nx^n=\sum_n(s_0+\dots+s_n)x^n=\frac{S(x)}{1-x}$. By the identity in the previous proof, $f(x)=(1-x)S(x)=(1-x)^2\sum_n(n+1)\sigma_nx^n$. The
weights $w_n(x)=(1-x)^2(n+1)x^n$ are non-negative, sum to $1$ since $\sum_n(n+1)x^n=(1-x)^{-2}$, and tend to $0$ for fixed $n$. Apply
Lemma~\ref{lem:regular} to $\sigma_n$.
\end{proof}

\begin{example}\label{ex:gregoryseries}
For $a_n=(-1)^n$ the partial sums are $1,0,1,0,\dots$, so $\sigma_n\to\frac12$ and the Abel limit is $f(x)=\frac1{1+x}\to\frac12$: both methods assign the
value Leibniz and Euler would have written down. Figure~\ref{fig:abelmeans} shows the three sequences. Abel summation is strictly stronger:
Exercise~\ref{ex:cesaroabel} gives a series that is Abel-summable and not Ces\`aro-summable.
\end{example}

\begin{figure}[ht]
\centering
\fig{ch10-abel}
\caption{The series $\sum(-1)^n$. The partial sums $s_n$ oscillate between $1$ and $0$ for ever; the Ces\`aro means $\sigma_n$ converge to $\frac12$;
and the Abel means $f(x)=\sum(-1)^nx^n$, drawn against $x$ on the same horizontal scale as $n=\frac x{1-x}$, approach $\frac12$ smoothly. Each method
is an average of the partial sums with weights that spread out as the parameter goes to its limit.}
\label{fig:abelmeans}
\end{figure}

\begin{think}
Every method above is an average of the partial sums, so it respects the ordinary sum. Could a summation method be \emph{strictly} stronger than all of
them at once and still be linear and respect ordinary sums? The Banach limits of Exercise~1.1 assign a value to every bounded sequence. What goes wrong if
you ask a method to sum $\sum_nn$ as well?
\end{think}

\section{Tauber's theorem}

Alfred Tauber published his theorem in 1897, in four pages of the \emph{Monatshefte}. He was not a professional researcher for most of his life: he
earned his living as an actuary and chief mathematician of the Ph\"onix insurance company in Vienna, lecturing at the university on the side. In June 1942,
aged seventy-six, he was deported to Theresienstadt, where he died.

\begin{theorem}[Tauber, 1897]\label{thm:tauber}
Let $f(x)=\sum_{n\ge0}a_nx^n$ for $|x|<1$. If $f(x)\to s$ as $x\to1^-$ and $na_n\to0$, then $\sum_na_n$ converges to $s$.
\end{theorem}
\begin{proof}
Let $\varepsilon_n=n|a_n|$, so $\varepsilon_n\to0$; put $\eta_N=\sup_{n>N}\varepsilon_n\to0$ and $\mu_N=\frac1N\sum_{n\le N}\varepsilon_n\to0$ (the averages of a null
sequence are null, by Lemma~\ref{lem:regular} applied to $\varepsilon_n$, or directly). Take $x=1-\frac1N$. Then
\[
s_N-f(x)=\sum_{n\le N}a_n(1-x^n)-\sum_{n>N}a_nx^n.
\]
For the first sum, $|1-x^n|=(1-x)(1+x+\dots+x^{n-1})\le n(1-x)=\frac nN$, so it is at most $\frac1N\sum_{n\le N}n|a_n|=\mu_N$. For the second,
$|a_n|\le\frac{\eta_N}n\le\frac{\eta_N}N$ for $n>N$, so it is at most $\frac{\eta_N}N\sum_{n>N}x^n\le\frac{\eta_N}N\cdot\frac1{1-x}=\eta_N$. Hence
$|s_N-f(1-\frac1N)|\le\mu_N+\eta_N\to0$, while $f(1-\frac1N)\to s$. So $s_N\to s$.
\end{proof}

\begin{remark}
In 1911 Littlewood proved the much harder theorem that $a_n=O(\frac1n)$ suffices: an Abel-summable series with terms $O(\frac1n)$ converges. Hardy
recalled that the proof cost Littlewood a great deal of work, and that no short proof was found for decades; Karamata's method below gives one of the
modern routes. We prove it nowhere in this book, and state it only for orientation.
\end{remark}

\begin{remark}
Tauber's theorem is sharp in the direction it points: the hypothesis cannot be dropped, since $\sum(-1)^n$ has $na_n\not\to0$. But it is not the right
hypothesis for arithmetic. In number theory the coefficients are non-negative and large, and the useful one-sided condition is positivity, which is what
Karamata and Wiener--Ikehara exploit below.
\end{remark}

\section{Karamata's theorem}

Hardy and Littlewood proved in 1914 that for $a_n\ge0$, an Abel-type asymptotic $\sum a_nx^n\sim\frac C{1-x}$ forces $\sum_{n\le N}a_n\sim CN$. Their proof
was long. In 1930 Jovan Karamata, then in Belgrade, published a proof of a page and a half whose whole engine is the Weierstrass approximation theorem. It
is one of the most elegant arguments in analysis, and we give it in the generality we need, for measures instead of sequences.

Let $\mu$ be a positive Borel measure on $[0,\infty[$ with distribution function $U(x)=\mu([0,x])$ and Laplace transform
\[
\omega(t)=\int_{[0,\infty[}e^{-tu}\dd\mu(u),\qq t>0,
\]
assumed finite for every $t>0$. We want to read the growth of $U$ at infinity off the behaviour of $\omega$ at $0$. Pure powers must be allowed a
correction factor that grows more slowly than any power.

\begin{definition}
A measurable function $\ell:\left]0,\infty\right[\to\left]0,\infty\right[$ is \emph{slowly varying} (at infinity) if $\frac{\ell(\lambda x)}{\ell(x)}\to1$ as
$x\to\infty$, for each fixed $\lambda>0$.
\end{definition}

Constants are slowly varying, and so are $\log x$, $(\log x)^\alpha$ and $\log\log x$; $x^\alpha$ is not, for $\alpha\ne0$. A function of the form $x^\rho\ell(x)$ with
$\ell$ slowly varying is called \emph{regularly varying of index $\rho$}: these are exactly the functions with a clean power-like growth, and they are the
natural output of a Tauberian theorem.

\begin{theorem}[Hardy, Littlewood, Karamata]\label{thm:karamata}
Let $\mu$, $U$ and $\omega$ be as above, $\rho\ge0$, and $\ell$ slowly varying. If
\[
\omega(t)\sim t^{-\rho}\,\ell\bigl(\tfrac1t\bigr)\qq\text{as }t\to0^+,
\]
then
\[
U(x)\sim\frac{x^\rho\ell(x)}{\Gamma(\rho+1)}\qq\text{as }x\to\infty.
\]
\end{theorem}
\begin{proof}
\emph{The functionals.} For a bounded Borel function $g$ on $[0,1]$ set
\[
I_t(g)=\frac1{\omega(t)}\int_{[0,\infty[}e^{-tu}g(e^{-tu})\dd\mu(u),\qq t>0.
\]
Each $I_t$ is linear, positive ($g\ge0$ implies $I_t(g)\ge0$) and normalised: $I_t(1)=1$. Consequently $|I_t(g)|\le\sup|g|$ and, for two functions,
$|I_t(g)-I_t(h)|\le\sup|g-h|$. Define the candidate limit
\[
I(g)=\frac1{\Gamma(\rho)}\int_0^\infty e^{-v}g(e^{-v})v^{\rho-1}\dd v\quad(\rho>0),\qq I(g)=g(1)\quad(\rho=0).
\]
Both are positive, linear and normalised, the first because $\int_0^\infty e^{-v}v^{\rho-1}\dd v=\Gamma(\rho)$.

\emph{Powers.} For an integer $k\ge0$ take $g(y)=y^k$. Then $e^{-tu}g(e^{-tu})=e^{-(k+1)tu}$, so $I_t(y^k)=\frac{\omega((k+1)t)}{\omega(t)}$. By hypothesis and
slow variation with $\lambda=\frac1{k+1}$,
\[
\frac{\omega((k+1)t)}{\omega(t)}\sim(k+1)^{-\rho}\frac{\ell\bigl(\frac1{(k+1)t}\bigr)}{\ell\bigl(\frac1t\bigr)}\longrightarrow(k+1)^{-\rho}.
\]
And $I(y^k)=\frac1{\Gamma(\rho)}\int_0^\infty e^{-(k+1)v}v^{\rho-1}\dd v=(k+1)^{-\rho}$ for $\rho>0$, while for $\rho=0$ both sides are $1$. So $I_t(g)\to I(g)$ for
$g$ a power, hence for every polynomial.

\emph{Continuous functions.} Let $g\in C([0,1])$ and $\varepsilon>0$. By the Weierstrass approximation theorem (Volume~6, Theorem~4.19) there is a polynomial $P$
with $\sup_{[0,1]}|g-P|<\varepsilon$. Then
\[
|I_t(g)-I(g)|\le|I_t(g)-I_t(P)|+|I_t(P)-I(P)|+|I(P)-I(g)|\le2\varepsilon+|I_t(P)-I(P)|,
\]
and the middle term tends to $0$. So $I_t(g)\to I(g)$ for every continuous $g$.

\emph{The indicator.} Let $g_0(y)=\frac1y\mathbf 1_{[e^{-1},1]}(y)$, so that $e^{-tu}g_0(e^{-tu})=\mathbf 1_{\{tu\le1\}}$ and
\[
I_t(g_0)=\frac{\mu([0,\frac1t])}{\omega(t)}=\frac{U(1/t)}{\omega(t)}.
\]
Also $I(g_0)=\frac1{\Gamma(\rho)}\int_0^1v^{\rho-1}\dd v=\frac1{\rho\Gamma(\rho)}=\frac1{\Gamma(\rho+1)}$ for $\rho>0$, and $I(g_0)=g_0(1)=1=\frac1{\Gamma(1)}$ for $\rho=0$. The function
$g_0$ is not continuous, so we squeeze it. For small $\delta>0$ let $\varphi_\delta$ be continuous, equal to $0$ on $[0,e^{-1}-\delta]$, to $1$ on $[e^{-1},1]$ and affine in
between, and let $\psi_\delta$ be continuous, $0$ on $[0,e^{-1}]$, $1$ on $[e^{-1}+\delta,1]$ and affine in between. Then $h^-_\delta=\frac1y\psi_\delta$ and $h^+_\delta=\frac1y\varphi_\delta$ are
continuous on $[0,1]$ (both vanish near $0$) and $h^-_\delta\le g_0\le h^+_\delta$, with
\[
0\le h^+_\delta-h^-_\delta\le\tfrac1y\mathbf 1_{J_\delta},\qq J_\delta=[e^{-1}-\delta,e^{-1}+\delta].
\]
By positivity, $I_t(h^-_\delta)\le I_t(g_0)\le I_t(h^+_\delta)$, so
\[
I(h^-_\delta)\le\liminf_{t\to0^+}I_t(g_0)\le\limsup_{t\to0^+}I_t(g_0)\le I(h^+_\delta),
\]
and it remains to check that $I(h^+_\delta)-I(h^-_\delta)\to0$ as $\delta\to0$. On $J_\delta$ we have $\frac1y\le2e$ for $\delta$ small, so the difference is at most
$2e\,I(\mathbf 1_{J_\delta})$. For $\rho>0$, $I(\mathbf 1_{J_\delta})=\frac1{\Gamma(\rho)}\int_{A_\delta}e^{-v}v^{\rho-1}\dd v$ where $A_\delta=\{v\ge0:e^{-v}\in J_\delta\}$ is an interval shrinking to
$\{1\}$, so the integral tends to $0$; for $\rho=0$, $\mathbf 1_{J_\delta}(1)=0$. Hence $I_t(g_0)\to\frac1{\Gamma(\rho+1)}$, that is
\[
U(1/t)\sim\frac{\omega(t)}{\Gamma(\rho+1)}\sim\frac{t^{-\rho}\ell(1/t)}{\Gamma(\rho+1)},
\]
which is the claim with $x=\frac1t$.
\end{proof}

Notice what the proof needed and what it did not. It needed positivity, to turn the convergence on polynomials into convergence on continuous functions by
a uniform estimate, and to squeeze a jump between continuous functions. It needed the hypothesis only along the fixed dilations $t\mapsto(k+1)t$, which is why
slow variation in its bare pointwise form is enough. It used no complex analysis at all.

\begin{corollary}[power series]\label{cor:karamataseries}
Let $a_n\ge0$ with $f(x)=\sum_{n\ge0}a_nx^n$ convergent for $0<x<1$, let $\rho\ge0$ and $\ell$ be slowly varying. If $f(e^{-t})\sim t^{-\rho}\ell(\frac1t)$ as
$t\to0^+$, then $\sum_{n\le x}a_n\sim\frac{x^\rho\ell(x)}{\Gamma(\rho+1)}$.
\end{corollary}
\begin{proof}
Apply Theorem~\ref{thm:karamata} to $\mu=\sum_{n\ge0}a_n\delta_n$: then $\omega(t)=\sum_na_ne^{-nt}=f(e^{-t})$ and $U(x)=\sum_{n\le x}a_n$.
\end{proof}

\begin{corollary}[a converse to Abel]\label{cor:abelconverse}
If $a_n\ge0$ and $f(x)=\sum a_nx^n\to s$ as $x\to1^-$, then $\sum_na_n$ converges, with sum $s$.
\end{corollary}
\begin{proof}
If $s=0$ then every $a_n=0$, since $f$ is non-decreasing on $[0,1[$ with non-negative coefficients. If $s>0$, apply
Corollary~\ref{cor:karamataseries} with $\rho=0$ and $\ell\equiv s$ (note $f(e^{-t})\to s$ as $t\to0^+$): $\sum_{n\le x}a_n\to s$.
\end{proof}

So for non-negative coefficients, Abel summability \emph{is} convergence. This is the first genuinely Tauberian statement of the chapter: a one-sided
condition on the terms, no smallness assumption, and a conclusion about the partial sums.

\section{Wiener's Tauberian theorem}

The methods so far all compare the transform with a single pure power. Wiener's theorem is of another kind. It answers the question: which averages of a
bounded function determine its behaviour at infinity? If $f\in L^1(\R)$ and $\phi\in L^\infty(\R)$, the convolution
\[
(\phi*f)(x)=\int_\R\phi(x-y)f(y)\dd y
\]
is a weighted average of $\phi$ around $x$ (Volume~6, Proposition~9.9(b): it is bounded and continuous). Suppose we know that this average tends to a limit as
$x\to\infty$. Does the \emph{same} hold for every other weight $g\in L^1$? Wiener's answer, from his hundred-page paper of 1932, is yes, provided the
Fourier transform of $f$ never vanishes --- and no, in general, if it vanishes somewhere. The algebraic heart of the matter was Chapter~16's theme: $L^1(\R)$
is a commutative Banach algebra under convolution, its characters are the points of the frequency line, and a function whose transform never vanishes lies
in no maximal ideal, so it generates everything. Here we prove it by hand, with two explicit kernels and a geometric series.

Throughout, $\hat g(\xi)=\int g(x)e^{-2\pi ix\xi}\dd x$, as in Chapter~18, and $\tau_yg(x)=g(x-y)$.

\emph{Two kernels.} Let
\[
F(x)=\Bigl(\frac{\sin\pi x}{\pi x}\Bigr)^2,\qq\Lambda(\xi)=(1-|\xi|)_+,
\]
so that $\hat F=\Lambda$: indeed $\Lambda=\mathbf 1_{[-\frac12,\frac12]}*\mathbf 1_{[-\frac12,\frac12]}$ has transform $\bigl(\frac{\sin\pi x}{\pi x}\bigr)^2$ by
Proposition~5.1(c), and $F$ and $\Lambda$ are even, so Fourier inversion (Theorem~5.2) exchanges them. For $a>0$ put $F_a(x)=aF(ax)$, whose transform is
$\Lambda(\frac\xi a)$. Note $F\ge0$ and $\int F=\hat F(0)=1$, so the family $(F_a)_{a>0}$ is an approximate identity in the sense of Volume~6: $F_a(x)=a F(ax)$ with
$F\ge0$ of integral $1$. Finally set
\[
V=2F_2-F_1,\qq \hat V(\xi)=2\Lambda\bigl(\tfrac\xi2\bigr)-\Lambda(\xi)=\begin{cases}1,&|\xi|\le1,\\2-|\xi|,&1\le|\xi|\le2,\\0,&|\xi|\ge2,\end{cases}
\]
a trapezoid: $\hat V=1$ on $[-1,1]$ and $\hat V=0$ outside $[-2,2]$. For $\delta>0$ write $V^{(\delta)}(x)=\delta V(\delta x)$, with transform $\hat V(\frac\xi\delta)$, equal to
$1$ on $[-\delta,\delta]$ and vanishing outside $[-2\delta,2\delta]$.

\begin{figure}[ht]
\centering
\fig{ch10-fejer}
\caption{Left: the kernel $k_\lambda(v)=\lambda\bigl(\frac{\sin\pi\lambda v}{\pi\lambda v}\bigr)^2$ for $\lambda=1$ and $\lambda=3$. It is non-negative, has
integral $1$, and concentrates at the origin as $\lambda$ grows. Right: its Fourier transform, the tent $(1-|\xi|/\lambda)_+$, which is $1$ at the origin and
has compact support. Non-negativity on one side and compact support on the other are exactly what the Tauberian arguments of this chapter need.}
\label{fig:fejerkernel}
\end{figure}

\begin{lemma}[local inverse]\label{lem:localinv}
Let $f\in L^1(\R)$ and $\xi_0\in\R$ with $\hat f(\xi_0)\ne0$. Then there are $h\in L^1(\R)$ and $\delta>0$ with $\hat h(\xi)\hat f(\xi)=1$ for $|\xi-\xi_0|\le\delta$.
\end{lemma}
\begin{proof}
\emph{Reduction to $\xi_0=0$, $\hat f(0)=1$.} Let $f_1(x)=\frac{f(x)e^{-2\pi i\xi_0x}}{\hat f(\xi_0)}$. The convention $\hat g(\xi)=\int g(x)e^{-2\pi ix\xi}\dd x$ gives $\hat f_1(\xi)=\frac{\hat f(\xi+\xi_0)}{\hat f(\xi_0)}$, and
$\hat f_1(0)=1$. If $h_1$ works for $f_1$ at $0$, set $h(x)=\frac{h_1(x)e^{+2\pi i\xi_0x}}{\hat f(\xi_0)}$. The same convention shifts the other way: $\hat h(\xi)=\frac{\hat h_1(\xi-\xi_0)}{\hat f(\xi_0)}$. Therefore
$\hat h(\xi)\hat f(\xi)=\hat h_1(\xi-\xi_0)\cdot\frac{\hat f(\xi)}{\hat f(\xi_0)}=\hat h_1(\xi-\xi_0)\hat f_1(\xi-\xi_0)$. On $|\xi-\xi_0|\le\delta$ the right side is $1$ whenever $\hat h_1\hat f_1=1$ on $[-\delta,\delta]$. So assume $\hat f(0)=\int f=1$.

\emph{A small remainder.} Let $w=f-V$, so $w\in L^1$ and $\int w=1-\hat V(0)=0$. Put $u_\delta=V^{(\delta)}*w$. Since $\int w=0$,
\[
u_\delta(x)=\int\bigl(V^{(\delta)}(x-y)-V^{(\delta)}(x)\bigr)w(y)\dd y,\qq\text{so}\qq\|u_\delta\|_1\le\int|w(y)|\,\bigl\|\tau_yV^{(\delta)}-V^{(\delta)}\bigr\|_1\dd y.
\]
A change of variable gives $\|\tau_yV^{(\delta)}-V^{(\delta)}\|_1=\|\tau_{\delta y}V-V\|_1$, which is at most $2\|V\|_1$ and tends to $0$ as $\delta\to0$ for each fixed $y$,
by continuity of translation in $L^1$ (Volume~6, Proposition~9.8). Dominated convergence gives $\|u_\delta\|_1\to0$. Fix $\delta<\frac12$ with $\|u\|_1<\frac12$, where
$u=u_\delta$.

\emph{Inversion.} Since $\delta<\frac12$, the transform $\hat V(\frac\xi\delta)$ is supported in $|\xi|\le2\delta<1$, where $\hat V=1$; hence
$\hat V(\frac\xi\delta)\hat V(\xi)=\hat V(\frac\xi\delta)$, that is $V^{(\delta)}*V=V^{(\delta)}$ (uniqueness of transforms, Theorem~5.2). Therefore
\[
V^{(\delta)}*f=V^{(\delta)}*V+V^{(\delta)}*w=V^{(\delta)}+u,
\]
and on $|\xi|\le\delta$, where $\hat V(\frac\xi\delta)=1$, this reads $\hat f(\xi)=1+\hat u(\xi)$. Now put
\[
h=V^{(\delta)}+\sum_{n\ge1}(-1)^nV^{(\delta)}*u^{*n},
\]
where $u^{*n}$ is the $n$-fold convolution. The series converges in $L^1$ because $\|V^{(\delta)}*u^{*n}\|_1\le\|V^{(\delta)}\|_1\|u\|_1^n\le\|V^{(\delta)}\|_1 2^{-n}$, so
$h\in L^1$, and taking transforms term by term,
\[
\hat h(\xi)=\hat V\bigl(\tfrac\xi\delta\bigr)\sum_{n\ge0}(-\hat u(\xi))^n=\frac{\hat V(\xi/\delta)}{1+\hat u(\xi)},
\]
the geometric series being legitimate since $|\hat u|\le\|u\|_1<\frac12$. For $|\xi|\le\delta$ this is $\frac1{1+\hat u(\xi)}=\frac1{\hat f(\xi)}$.
\end{proof}

\begin{theorem}[Wiener, 1932]\label{thm:wienerl1}
Let $f\in L^1(\R)$ with $\hat f(\xi)\ne0$ for every $\xi\in\R$. Then the closed linear span of the translates $\tau_yf$, $y\in\R$, is the whole of $L^1(\R)$.
\end{theorem}
\begin{proof}
Let $J$ be that closed span.

\emph{Step 1: $g*f\in J$ for every $g\in L^1$.} The map $y\mapsto\tau_yf$ is continuous from $\R$ to $L^1$ (Volume~6, Proposition~9.8), hence uniformly
continuous on a compact interval, so for $a<b$ the Riemann sums of $\int_a^b\tau_yf\dd y$ converge in $L^1$; the limit lies in $J$, and it is $\mathbf
1_{[a,b]}*f$. By linearity $s*f\in J$ for every step function $s$. Step functions are dense in $L^1$ (Volume~6, Chapter~7), and
$\|s*f-g*f\|_1\le\|s-g\|_1\|f\|_1$, so $g*f\in J$.

\emph{Step 2: if $\hat g$ has compact support, then $g=h*f$ for some $h\in L^1$.} Let $K=\operatorname{supp}\hat g$, a compact set. For each $\xi\in K$,
Lemma~\ref{lem:localinv} gives $h_\xi\in L^1$ and an open interval $I_\xi\ni\xi$ with $\hat h_\xi\hat f=1$ on $I_\xi$. Finitely many $I_{\xi_1},\dots,I_{\xi_N}$
cover $K$, and by compactness there is $\rho>0$ such that every point of $K$ has its $\rho$-neighbourhood inside one of them. (Otherwise there are
$x_j\in K$ whose $\frac1j$-neighbourhoods meet the complement of every $I_{\xi_i}$; a convergent subsequence $x_j\to x\in K$ then has $x$ outside the open
union, which is absurd.)

Choose $\eta$ with $2\eta<\rho$ and consider the hat functions $\Lambda_k(\xi)=\Lambda(\frac\xi\eta-k)$ for $k\in\Z$: piecewise affine, supported on
$[(k-1)\eta,(k+1)\eta]$, and $\sum_{k\in\Z}\Lambda_k=1$ everywhere, since on each interval $[k\eta,(k+1)\eta]$ only the two hats $k$ and $k+1$ are non-zero and they
add to $1$. Each $\Lambda_k$ is a Fourier transform: $\Lambda_k=\hat p_k$ with $p_k(x)=\eta F(\eta x)e^{2\pi ik\eta x}\in L^1$. Let $S$ be the (finite) set of $k$ for which
$\operatorname{supp}\Lambda_k$ meets $K$; for such $k$ that support has width $2\eta<\rho$ and meets $K$, so it lies in some $I_{\xi_{i(k)}}$. Since $\hat g$ vanishes
off $K$ and $\sum_{k\in S}\Lambda_k=1$ on $K$, we get $\hat g=\sum_{k\in S}\hat g\Lambda_k$, that is $g=\sum_{k\in S}g*p_k$. For $k\in S$ the function $\hat g\Lambda_k$ is
supported in $I_{\xi_{i(k)}}$, where $\hat h_{\xi_{i(k)}}\hat f=1$, so
\[
\hat g\Lambda_k=\hat g\Lambda_k\hat h_{\xi_{i(k)}}\hat f,\qq\text{i.e.}\qq g*p_k=\bigl(g*p_k*h_{\xi_{i(k)}}\bigr)*f.
\]
Summing over $k\in S$ gives $g=h*f$ with $h=\sum_{k\in S}g*p_k*h_{\xi_{i(k)}}\in L^1$.

\emph{Step 3: density.} For $g\in L^1$, the transform of $F_a*g$ is $\Lambda(\frac\xi a)\hat g$, supported in $[-a,a]$, and $\|F_a*g-g\|_1\to0$ as $a\to\infty$
because $(F_a)$ is an approximate identity (Volume~6, Theorem~9.10(a)). So the functions with compactly supported transform are dense in $L^1$. By Steps 1
and 2 they all lie in $J$, and $J$ is closed, so $J=L^1(\R)$.
\end{proof}

\begin{theorem}[Wiener's Tauberian theorem]\label{thm:wienertaub}
Let $f\in L^1(\R)$ with $\hat f$ nowhere zero, let $\phi\in L^\infty(\R)$ and $A\in\C$, and suppose
\[
(\phi*f)(x)\longrightarrow A\int_\R f\qq\text{as }x\to+\infty.
\]
Then $(\phi*g)(x)\to A\int_\R g$ as $x\to+\infty$, for every $g\in L^1(\R)$.
\end{theorem}
\begin{proof}
Let $G=\{g\in L^1:(\phi*g)(x)\to A\int g\text{ as }x\to+\infty\}$. It is a linear subspace containing $f$. It is translation-invariant, since
$\phi*\tau_yg=\tau_y(\phi*g)$ has the same limit and $\int\tau_yg=\int g$. It is closed: if $g_j\to g$ in $L^1$ with $g_j\in G$, then given $\varepsilon>0$ pick $j$
with $\|g-g_j\|_1<\varepsilon$; since $\|\phi*g-\phi*g_j\|_\infty\le\|\phi\|_\infty\varepsilon$ and $|\int g-\int g_j|\le\varepsilon$, the limit superior of
$|(\phi*g)(x)-A\int g|$ is at most $(\|\phi\|_\infty+|A|)\varepsilon$ for every $\varepsilon$. So $G$ is a closed translation-invariant subspace containing $f$, hence
contains the closed span of the translates of $f$, which is $L^1(\R)$ by Theorem~\ref{thm:wienerl1}.
\end{proof}

\begin{remark}
The hypothesis is sharp in a strong sense: if $\hat f(\xi_0)=0$, take $\phi(x)=e^{2\pi i\xi_0x}$. Then $(\phi*f)(x)=e^{2\pi i\xi_0x}\hat f(\xi_0)=0$ for all $x$, while
$(\phi*g)(x)=e^{2\pi i\xi_0x}\hat g(\xi_0)$ does not tend to $0$ for any $g$ with $\hat g(\xi_0)\ne0$. One vanishing frequency, and the averages against $f$ are
blind to a whole oscillation of $\phi$. Exercise~\ref{ex:wienerZ} does the same story on $\Z$, where Chapter~16 already gave us the key.
\end{remark}

\begin{beauty}
Wiener's proof is the first clear instance of a \emph{local-to-global} principle in analysis: Lemma~\ref{lem:localinv} inverts $\hat f$ near one frequency, a
partition of unity glues finitely many such inversions, and an approximate identity reduces the general case to the compactly supported one. There is no
global formula for $h$; nothing like $\frac1{\hat f}$ belongs to the algebra, since $\frac1{\hat f}$ need not tend to $0$ at infinity and so need not be a
transform at all. The theorem is true because a frequency is a local notion, and $L^1$ can be cut up and reassembled along frequencies.
\end{beauty}

\section{Wiener--Ikehara and the prime number theorem}

Shikao Ikehara was Wiener's student at MIT. In 1931 he published an extension of a theorem of Landau that has since become the standard analytic bridge
from a Dirichlet series to its coefficients: a single pole at $s=1$, continuity elsewhere on the line, non-negative coefficients, and the partial sums are
asymptotically linear. Ikehara deduced it from Wiener's theorem of the previous section; Bochner then noticed that the Fej\'er kernel does the work
directly, and that is the proof below. We use the kernel
\[
k_\lambda(v)=\lambda\Bigl(\frac{\sin\pi\lambda v}{\pi\lambda v}\Bigr)^2=F_\lambda(v),\qq\hat k_\lambda(\xi)=\Bigl(1-\frac{|\xi|}\lambda\Bigr)_+,
\]
of Figure~\ref{fig:fejerkernel}: non-negative, of integral $1$, with transform supported in $[-\lambda,\lambda]$.

\begin{theorem}[Wiener--Ikehara]\label{thm:ikehara}
Let $a_n\ge0$ be such that $F(s)=\sum_{n\ge1}a_nn^{-s}$ converges for $\operatorname{Re}s>1$, and suppose that for some $C\ge0$ the function
$F(s)-\frac C{s-1}$ extends to a continuous function on the closed half-plane $\operatorname{Re}s\ge1$. Then
\[
A(x)=\sum_{n\le x}a_n\sim Cx\qq\text{as }x\to\infty.
\]
\end{theorem}
\begin{proof}
\emph{Setting up.} For $\operatorname{Re}s>1$, summing by parts (or integrating $\int_n^\infty x^{-s-1}\dd x$ term by term),
\[
\frac{F(s)}s=\int_1^\infty A(x)x^{-s-1}\dd x=\int_0^\infty S(v)e^{-sv}\dd v,\qq S(v)=A(e^v),
\]
where $S$ is non-decreasing and right-continuous. Let $\psi(v)=e^{-v}S(v)\ge0$; the claim is $\psi(v)\to C$. Set
\[
H(s)=\frac{F(s)}s-\frac C{s-1}.
\]
Since $\frac C{s(s-1)}+\frac Cs=\frac{C(1+(s-1))}{s(s-1)}=\frac C{s-1}$, we may rewrite $H$ as
\[
H(s)=\frac1s\Bigl(F(s)-\frac C{s-1}\Bigr)-\frac Cs.
\]
Both terms on the right are continuous on $\operatorname{Re}s\ge1$ by hypothesis (and $s\ne0$ there), so $H$ is
continuous there, and holomorphic on $\operatorname{Re}s>1$. For $\varepsilon>0$ put $\psi_\varepsilon(v)=(\psi(v)-C)e^{-\varepsilon v}$. It lies in $L^1(0,\infty)$, since
$\int_0^\infty\psi(v)e^{-\varepsilon v}\dd v=\frac{F(1+\varepsilon)}{1+\varepsilon}<\infty$, and, computing both Laplace integrals,
\[
\int_0^\infty\psi_\varepsilon(v)e^{-2\pi i\xi v}\dd v=\frac{F(s)}s-\frac C{\varepsilon+2\pi i\xi}=H(s),\qq s=1+\varepsilon+2\pi i\xi.
\]

\emph{Smoothing.} Multiply by the tent $\bigl(1-\frac{|\xi|}\lambda\bigr)_+$, which is $\hat k_\lambda$, and integrate against $e^{2\pi i\xi v}$. By Fubini
(legitimate: $\psi_\varepsilon\in L^1$ and the $\xi$-range is bounded) and Fourier inversion for $k_\lambda$,
\[
\int_{-\lambda}^\lambda H(1+\varepsilon+2\pi i\xi)\Bigl(1-\frac{|\xi|}\lambda\Bigr)e^{2\pi i\xi v}\dd\xi=\int_0^\infty\psi_\varepsilon(w)k_\lambda(v-w)\dd w.\tag{$*$}\label{eq:ikesmooth}
\]

\emph{$\psi$ is bounded.} Fix $\lambda=1$. The left side of \eqref{eq:ikesmooth} is at most $2M$ in absolute value, where $M$ is the supremum of $|H|$ on the
compact set $\{1+\varepsilon+2\pi i\xi:0\le\varepsilon\le1,|\xi|\le1\}$, finite by continuity. Since $k_1\ge0$ and $C\int_0^\infty e^{-\varepsilon w}k_1(v-w)\dd w\le C$, the right side gives
\[
\int_0^\infty\psi(w)e^{-\varepsilon w}k_1(v-w)\dd w\le2M+C\qq(0<\varepsilon\le1).
\]
For $w\in[v,v+\frac12]$ we have $|v-w|\le\frac12$, so $\frac{\sin\pi(v-w)}{\pi(v-w)}\ge\frac2\pi$ and $k_1(v-w)\ge\frac4{\pi^2}$; and $S(w)\ge S(v)$ gives
$\psi(w)\ge e^{-(w-v)}\psi(v)\ge e^{-1/2}\psi(v)$. Letting $\varepsilon\to0$ in the integral over this bounded interval, where $\psi$ is bounded,
\[
\frac4{\pi^2}\cdot\frac12\cdot e^{-1/2}\psi(v)\le2M+C,
\]
so $\psi$ is bounded, say by $B$. (In particular $A(x)=O(x)$.)

\emph{Letting $\varepsilon\to0$.} Now fix $\lambda>0$ and $v$. On the left of \eqref{eq:ikesmooth}, $H(1+\varepsilon+2\pi i\xi)\to H(1+2\pi i\xi)$ uniformly for $|\xi|\le\lambda$,
because $H$ is uniformly continuous on the compact set $\{1+\varepsilon+2\pi i\xi:0\le\varepsilon\le1,|\xi|\le\lambda\}$. On the right, $|\psi_\varepsilon|\le(B+C)$ and
$k_\lambda(v-\cdot)\in L^1$, so dominated convergence applies. Hence
\[
\int_{-\lambda}^\lambda H(1+2\pi i\xi)\Bigl(1-\frac{|\xi|}\lambda\Bigr)e^{2\pi i\xi v}\dd\xi=\int_0^\infty(\psi(w)-C)k_\lambda(v-w)\dd w.
\]

\emph{Letting $v\to\infty$.} The function $\xi\mapsto H(1+2\pi i\xi)(1-\frac{|\xi|}\lambda)_+$ is continuous with compact support, hence in $L^1(\R)$, so the left
side tends to $0$ as $v\to\infty$ by the Riemann--Lebesgue lemma (Proposition~5.1(b)). On the right, substitute $w=v+r$ and use $|\psi-C|\le B+C$ together
with $\int_{|r|>v}k_\lambda\to0$ to discard the range $r<-v$. We obtain, for every $\lambda>0$,
\[
\lim_{v\to\infty}\int_\R\bigl(\psi(v+r)-C\bigr)k_\lambda(r)\dd r=0.\tag{$**$}\label{eq:ikelimit}
\]

\emph{De-smoothing.} Let $\psi^*=\limsup_{v\to\infty}\psi(v)$ and $\psi_*=\liminf_{v\to\infty}\psi(v)$, both finite. Fix $\delta>0$ and write
$m=m(\delta,\lambda)=\int_{-\delta}^\delta k_\lambda$, which tends to $1$ as $\lambda\to\infty$ since $k_\lambda$ is an approximate identity of integral $1$.

Take $v_j\to\infty$ with $\psi(v_j)\to\psi^*$ and apply \eqref{eq:ikelimit} at the points $v_j+\delta$. For $|r|\le\delta$ we have $v_j+\delta+r\ge v_j$, so
monotonicity of $S$ gives $\psi(v_j+\delta+r)\ge e^{-(\delta+r)}\psi(v_j)\ge e^{-2\delta}\psi(v_j)$; elsewhere $\psi-C\ge-C$. Hence
\[
0=\lim_j\int_\R\bigl(\psi(v_j+\delta+r)-C\bigr)k_\lambda(r)\dd r\ge\bigl(e^{-2\delta}\psi^*-C\bigr)m-C(1-m).
\]
Letting $\lambda\to\infty$, so that $m\to1$, gives $e^{-2\delta}\psi^*\le C$, and then $\delta\to0$ gives $\psi^*\le C$.

Take $v_j\to\infty$ with $\psi(v_j)\to\psi_*$ and apply \eqref{eq:ikelimit} at the points $v_j-\delta$. For $|r|\le\delta$ we have $v_j-\delta+r\le v_j$, so
$\psi(v_j-\delta+r)\le e^{\delta-r}\psi(v_j)\le e^{2\delta}\psi(v_j)$; elsewhere $|\psi-C|\le B+C$. Hence
\[
0\le\bigl(e^{2\delta}\psi_*-C\bigr)m+(B+C)(1-m),
\]
and letting $\lambda\to\infty$ then $\delta\to0$ gives $\psi_*\ge C$. So $\psi(v)\to C$, that is $A(e^v)\sim Ce^v$.
\end{proof}

\begin{think}
The hypothesis is continuity on the \emph{whole} line $\operatorname{Re}s=1$, not just at $s=1$. Exercise~\ref{ex:ikeharaline} builds non-negative $a_n$ whose
Dirichlet series has, besides the pole at $1$, two poles at $1\pm i$, and whose partial sums oscillate for ever between $0.29x$ and $1.71x$. Where in the
proof above would the argument break for that example? And notice what the theorem does \emph{not} give: no rate. Why can no theorem with these hypotheses
give one?
\end{think}

\begin{corollary}[the prime number theorem]\label{cor:pntikehara}
$\psi(x)=\sum_{n\le x}\Lambda(n)\sim x$, and hence $\pi(x)\sim\frac x{\log x}$.
\end{corollary}
\begin{proof}
Take $a_n=\Lambda(n)\ge0$, so $F(s)=-\frac{\zeta'}\zeta(s)$ for $\operatorname{Re}s>1$ (Volume~4, Theorem~10.3 and its logarithmic derivative). The function $\zeta$ is
holomorphic on $\operatorname{Re}s>0$ except for a simple pole at $s=1$ (Volume~4, Proposition~10.6), has no zero in $\operatorname{Re}s>1$ by the Euler product, and
none on the line $\operatorname{Re}s=1$ by Theorem~5.2 of Volume~5. So $-\frac{\zeta'}\zeta$ is holomorphic on an open set containing $\{\operatorname{Re}s\ge1\}$ except at
$s=1$, where $\zeta(s)=\frac1{s-1}+O(1)$ makes $-\frac{\zeta'}\zeta$ a simple pole of residue $1$; hence $F(s)-\frac1{s-1}$ is holomorphic near $s=1$ too, and in
particular continuous on $\operatorname{Re}s\ge1$. Theorem~\ref{thm:ikehara} with $C=1$ gives $\psi(x)\sim x$, and the passage to $\pi(x)$ is Volume~5,
Theorem~5.5.
\end{proof}

\begin{corollary}[progressions]\label{cor:pntapikehara}
Let $q\ge1$ and $a$ be prime to $q$. Then $\psi(x;q,a)=\sum_{n\le x,\,n\equiv a\,(q)}\Lambda(n)\sim\frac x{\varphi(q)}$.
\end{corollary}
\begin{proof}
By orthogonality of characters, $a_n=\Lambda(n)\mathbf 1_{n\equiv a\,(q)}\ge0$ has Dirichlet series
\[
F(s)=\frac1{\varphi(q)}\sum_\chi\overline\chi(a)\Bigl(-\frac{L'}L(s,\chi)\Bigr).
\]
For $\chi\ne\chi_0$, $L(\cdot,\chi)$ is holomorphic on $\operatorname{Re}s>0$ and non-zero on $\operatorname{Re}s\ge1$ (Volume~5, Theorem~5.2, which covers $s=1$), so its
logarithmic derivative is continuous on the closed half-plane. For $\chi_0$, $L(s,\chi_0)=\zeta(s)\prod_{p\mid q}(1-p^{-s})$, so
$-\frac{L'}L(s,\chi_0)=-\frac{\zeta'}\zeta(s)-\sum_{p\mid q}\frac{\log p}{p^s-1}$, whose last term is holomorphic for $\operatorname{Re}s>0$. Hence
$F(s)-\frac1{\varphi(q)}\frac1{s-1}$ is continuous on $\operatorname{Re}s\ge1$, and Theorem~\ref{thm:ikehara} applies with $C=\frac1{\varphi(q)}$.
\end{proof}

Volume~5 proved both corollaries by Newman's contour argument, which is shorter still but needs the same input: $\zeta$ and the $L(\cdot,\chi)$ must not
vanish on the line. The point of the Tauberian route is that it is a general theorem about Dirichlet series with non-negative coefficients, and we now use
it on a problem where no contour integral is in sight.

\section{Counting sums of two squares}

Call $n\ge1$ \emph{representable} if $n=x^2+y^2$ for integers $x,y$, and let $b(n)=1$ if $n$ is representable and $0$ otherwise. Volume~1, Theorem~7.3 says
that $b(n)=1$ exactly when every prime $q\equiv3\pmod4$ divides $n$ to an even power. How many representable integers are there up to $x$? In 1908 Landau
proved
\[
B(x)=\sum_{n\le x}b(n)\sim\frac{Kx}{\sqrt{\log x}},\qq K=0.76422365\ldots,
\]
and in 1913, in his first letter to Hardy, Ramanujan announced the same result (with an error term he could not justify). The constant $K$ is now named
after both of them. Landau's proof needs a careful contour analysis around a branch point, because $x/\sqrt{\log x}$ is not a nice power. Karamata's theorem
gives the logarithmically weighted version for free, and the weighted version already contains the constant.

\begin{proposition}\label{prop:twosqeuler}
For $\operatorname{Re}s>1$ the series $F(s)=\sum_{n\ge1}b(n)n^{-s}$ converges absolutely and
\[
F(s)=\frac1{1-2^{-s}}\prod_{p\equiv1(4)}\frac1{1-p^{-s}}\prod_{q\equiv3(4)}\frac1{1-q^{-2s}}.
\]
Moreover
\[
F(s)^2=\zeta(s)L(s,\chi_4)\cdot\frac1{1-2^{-s}}\prod_{q\equiv3(4)}\frac1{1-q^{-2s}},
\]
and
\[
F(\sigma)\sim\frac{K\sqrt\pi}{\sqrt{\sigma-1}}\qq\text{as }\sigma\to1^+,\qq K=\frac1{\sqrt2}\prod_{q\equiv3(4)}\Bigl(1-\frac1{q^2}\Bigr)^{-1/2}.
\]
\end{proposition}
\begin{proof}
By Volume~1, Theorem~7.3, $b$ is the indicator of a set of integers described by independent conditions on each prime: all exponents are allowed at $2$ and
at $p\equiv1$, only even exponents at $q\equiv3$. So $b$ is multiplicative, and the Euler product (Volume~4, Theorem~10.3, whose proof only uses multiplicativity
and absolute convergence) gives, for $\operatorname{Re}s>1$,
\[
F(s)=\prod_p\sum_{k\ge0}b(p^k)p^{-ks},
\]
where the local factor is $\sum_kp^{-ks}=(1-p^{-s})^{-1}$ for $p=2$ and $p\equiv1$, and $\sum_{k\text{ even}}q^{-ks}=(1-q^{-2s})^{-1}$ for $q\equiv3$. This is the
first formula. Next, with $\chi_4$ the non-trivial character mod $4$ ($\chi_4(p)=1$ for $p\equiv1$, $-1$ for $p\equiv3$),
\[
\zeta(s)L(s,\chi_4)=\frac1{1-2^{-s}}\prod_{p\equiv1(4)}\frac1{(1-p^{-s})^2}\prod_{q\equiv3(4)}\frac1{1-q^{-2s}},
\]
since $(1-q^{-s})^{-1}(1+q^{-s})^{-1}=(1-q^{-2s})^{-1}$. Squaring the first formula and comparing gives the second. Finally let $\sigma\to1^+$ along the reals. The
product $P(s)=\prod_{q\equiv3}(1-q^{-2s})^{-1}$ converges absolutely and is continuous for $\operatorname{Re}s>\frac12$, with $P(1)=\prod_{q\equiv3}(1-q^{-2})^{-1}$;
also $(\sigma-1)\zeta(\sigma)\to1$ and $L(\sigma,\chi_4)\to L(1,\chi_4)=\frac\pi4$ by continuity (Volume~5, Chapter~3, Example~3.11). Hence
\[
(\sigma-1)F(\sigma)^2\longrightarrow\frac\pi4\cdot2\cdot P(1)=\frac\pi2P(1)=\pi K^2,
\]
and $F(\sigma)>0$, so $F(\sigma)\sqrt{\sigma-1}\to K\sqrt\pi$.
\end{proof}

\begin{theorem}\label{thm:twosqlog}
As $x\to\infty$,
\[
\sum_{\substack{n\le x\\ n=\square+\square}}\frac1n\sim2K\sqrt{\log x},\qq K=0.7642236535\ldots
\]
\end{theorem}
\begin{proof}
Let $\mu=\sum_{n\ge1}\frac{b(n)}n\delta_{\log n}$, a positive measure on $[0,\infty[$. Its distribution function is
$U(y)=\sum_{n\le e^y}\frac{b(n)}n$, and its Laplace transform is
\[
\omega(t)=\sum_{n\ge1}\frac{b(n)}ne^{-t\log n}=\sum_{n\ge1}b(n)n^{-1-t}=F(1+t),
\]
finite for every $t>0$. By Proposition~\ref{prop:twosqeuler}, $\omega(t)\sim K\sqrt\pi\,t^{-1/2}$, which is the hypothesis of Theorem~\ref{thm:karamata} with
$\rho=\frac12$ and $\ell\equiv K\sqrt\pi$. Therefore
\[
U(y)\sim\frac{K\sqrt\pi\,y^{1/2}}{\Gamma(\frac32)}=\frac{K\sqrt\pi\sqrt y}{\frac12\sqrt\pi}=2K\sqrt y.
\]
Put $y=\log x$.
\end{proof}

\begin{remark}
The unweighted count is harder and we do not prove it: $B(x)\sim\frac{Kx}{\sqrt{\log x}}$ (Landau, 1908). The two statements are consistent, since
$\int_2^x\frac{\dd t}{t\sqrt{\log t}}=2\sqrt{\log x}+O(1)$, but neither follows from the other by elementary means: a Tauberian theorem of the kind we proved
cannot see the difference between $B(x)\sim\frac{Kx}{\sqrt{\log x}}$ and a count that oscillates around it by a factor $1+o(1)$ in the logarithmic average
only. What Karamata's theorem buys is the exact constant with no contour integral; what it costs is the logarithmic weight. Delange's 1954 generalisation of
Ikehara's theorem handles the branch point directly and gives Landau's form, together with a whole asymptotic expansion.
\end{remark}

\begin{beauty}
The constant $K$ is built from the primes $q\equiv3\pmod4$ alone: $K=\frac1{\sqrt2}\prod_{q\equiv3(4)}(1-q^{-2})^{-1/2}$. It says that the density of
representable integers, measured on the logarithmic scale, is decided entirely by how often the ``forbidden'' primes appear to an odd power. The square
root in $\sqrt{\log x}$ comes from the fact that the representable integers are, in multiplicative terms, half of the primes' worth of freedom: the
Dirichlet series has a square-root singularity because $F^2$, and not $F$, is the one with a pole. The same mechanism produces a factor $(\log x)^{\alpha-1}$
whenever a multiplicative condition leaves a proportion $\alpha$ of the primes free, which is Landau's theorem in its general form and the reason that
constants like $K$, built from one congruence class of primes, appear all over multiplicative number theory.
\end{beauty}

\section{Exercises}

\begin{exercise}[$\star$]\label{ex:cesaroabel}
\emph{Abel beats Ces\`aro.} Let $a_n=(-1)^n(n+1)$. (a) Show that $f(x)=\frac1{(1+x)^2}$, so the series is Abel-summable with sum $\frac14$. (b) Compute
$s_n$ and $\sigma_n$ and show that $\sigma_n$ has two limit points, so the series is not Ces\`aro-summable. (c) Why does Tauber's theorem not apply?
\end{exercise}

\begin{exercise}[$\star$]\label{ex:divisorkaramata}
\emph{Divisors, by Karamata.} Exercise~9.7 showed $\sum_{n\ge1}d(n)e^{-nx}\sim\frac{\log\frac1x}x$ as $x\to0^+$. Deduce $\sum_{n\le N}d(n)\sim N\log N$.
(Dirichlet's elementary proof gives the much sharper $N\log N+(2\gamma-1)N+O(\sqrt N)$; the point is how little the Tauberian route needs to know.)
\end{exercise}

\begin{exercise}[$\star\star$]\label{ex:partitionsup}
\emph{An upper bound for partitions.} Let $p(n)$ be the number of partitions of $n$ and $P(z)=\prod_{n\ge1}(1-z^n)^{-1}=\sum_np(n)z^n$ for $|z|<1$.
(a) Show that $\log P(e^{-t})=\sum_{m\ge1}\frac1m\frac1{e^{mt}-1}$ and deduce $\log P(e^{-t})\le\frac{\zeta(2)}t=\frac{\pi^2}{6t}$. (b) Show
$t\log P(e^{-t})\to\zeta(2)$ as $t\to0^+$. (c) Deduce $\log p(n)\le\pi\sqrt{\frac{2n}3}$ for every $n\ge1$. (The true asymptotic, Hardy and Ramanujan, 1918, is
$p(n)\sim\frac1{4n\sqrt3}e^{\pi\sqrt{2n/3}}$: the exponential is exactly the bound in (c), so the one-line argument loses only a polynomial factor. The
matching lower bound needs a different idea, and the constant $\frac1{4n\sqrt3}$ needs the circle method.)
\end{exercise}

\begin{exercise}[$\star\star$]\label{ex:quicksort}
\emph{Quicksort, by Karamata.} Let $C_n$ be the average number of comparisons that quicksort makes on $n$ distinct items, with the pivot uniform:
$C_0=0$ and $C_n=n-1+\frac2n\sum_{k=0}^{n-1}C_k$ for $n\ge1$. (a) Show that $C(z)=\sum_nC_nz^n$ satisfies $C(z)=\frac{2}{(1-z)^2}\bigl(\log\frac1{1-z}-z\bigr)$.
(b) Show that $(C_n)$ is non-decreasing. (c) Apply Corollary~\ref{cor:karamataseries} to get $\sum_{k\le n}C_k\sim n^2\log n$, and deduce from (b), by
comparing $\sum_{k\le n}$ with $\sum_{k\le\theta n}$, that $C_n\sim2n\log n$.
\end{exercise}

\begin{exercise}[$\star$]\label{ex:wienerZ}
\emph{Wiener on the integers.} Let $a\in\ell^1(\Z)$ with $\hat a(\theta)=\sum_ka_ke^{ik\theta}\ne0$ for all $\theta$, and let $x\in\ell^\infty(\Z)$. Show that if
$(a*x)_n\to s$ as $n\to\infty$, then $x_n\to\frac s{\hat a(0)}$. (Use Theorem~3.11.) Deduce: if $\frac{2x_n+x_{n-1}}3\to c$ then $x_n\to c$, whereas
$\frac{x_n+x_{n-1}}2\to0$ does not force $x_n$ to converge.
\end{exercise}

\begin{exercise}[$\star\star$]\label{ex:ikeharaline}
\emph{One pole off the real axis is enough to spoil everything.} Let $a_n=1+\cos(\log n)\ge0$. (a) Show that
$\sum_na_nn^{-s}=\zeta(s)+\frac12\bigl(\zeta(s-i)+\zeta(s+i)\bigr)$ for $\operatorname{Re}s>1$. (b) Show that
$\sum_{n\le x}a_n=x+\frac x{\sqrt2}\cos\bigl(\log x-\frac\pi4\bigr)+O(\log x)$, so that $\frac1x\sum_{n\le x}a_n$ oscillates for ever between $1-\frac1{\sqrt2}$
and $1+\frac1{\sqrt2}$. (c) Which hypothesis of Theorem~\ref{thm:ikehara} fails?
\end{exercise}

\section{Hints}

\begin{hint}{ex:cesaroabel}
Differentiate the geometric series. The partial sums of $(s_n)$ are $1,0,2,0,3,0,\dots$
\end{hint}
\begin{hint}{ex:divisorkaramata}
$\log\frac1x$ is slowly varying in $\frac1x$, so $\rho=1$ and $\ell(y)=\log y$.
\end{hint}
\begin{hint}{ex:partitionsup}
$\frac1{e^u-1}\le\frac1u$. For (b) write $\frac1{e^{mt}-1}=\frac1{mt}g(mt)$ with $g(u)=\frac u{e^u-1}\le1$ and use dominated convergence on the sum over $m$. For
(c), $p(n)z^n\le P(z)$, and optimise $t$.
\end{hint}
\begin{hint}{ex:quicksort}
For (a) use $\sum_nH_nz^n=\frac1{1-z}\log\frac1{1-z}$ and the closed form $C_n=2(n+1)H_n-4n$, or turn the recurrence into a differential equation for $C$. For
(b) induct, using that the recurrence expresses $C_n$ in terms of the earlier ones.
\end{hint}
\begin{hint}{ex:wienerZ}
Wiener's lemma gives $b\in\ell^1$ with $a*b=\delta_0$; then $x=b*(a*x)$ and dominated convergence does the rest.
\end{hint}
\begin{hint}{ex:ikeharaline}
$\cos(\log n)=\frac{n^i+n^{-i}}2$. For (b) compare the sum with $\operatorname{Re}\int_1^xt^i\dd t$.
\end{hint}

\section{Solutions}

\begin{solution}{ex:cesaroabel}
\idea Abel's weights $(1-x)x^n$ damp the tail far more strongly than Ces\`aro's $\frac1{n+1}$, so they can sum wilder series.
\writeup (a) $\sum_n(-1)^n(n+1)x^n=\sum_n(n+1)(-x)^n=\frac1{(1-(-x))^2}=\frac1{(1+x)^2}\to\frac14$. (b) $s_{2m}=m+1$ and $s_{2m+1}=-(m+1)$: indeed
$s_{2m+1}=s_{2m}-(2m+2)=m+1-2m-2$. The partial sums of $(s_n)$ are therefore $1,0,2,0,3,0,\dots$, so $(n+1)\sigma_n=m+1$ for $n=2m$ and $0$ for $n=2m+1$: thus
$\sigma_{2m}=\frac{m+1}{2m+1}\to\frac12$ and $\sigma_{2m+1}=0$. Two limit points, so no Ces\`aro sum. (c) $na_n=(-1)^nn(n+1)$ does not tend to $0$; nor are the $a_n$
non-negative, so Corollary~\ref{cor:abelconverse} does not apply either. Nothing is wrong: the series simply does not converge, and Abel's method is
genuinely stronger than convergence.
\end{solution}

\begin{solution}{ex:divisorkaramata}
\idea The generating function has a pole of order $1$ with a logarithmic correction, which is exactly what a slowly varying $\ell$ is for.
\writeup With $a_n=d(n)\ge0$ and $x=t$, Exercise~9.7 gives $\sum_nd(n)e^{-nt}\sim t^{-1}\log\frac1t=t^{-1}\ell(\frac1t)$ with $\ell(y)=\log y$, which is slowly
varying: $\frac{\log\lambda y}{\log y}\to1$. By Corollary~\ref{cor:karamataseries} with $\rho=1$, $\sum_{n\le N}d(n)\sim\frac{N\log N}{\Gamma(2)}=N\log N$.
\end{solution}

\begin{solution}{ex:partitionsup}
\idea An upper bound for a coefficient of a power series with non-negative coefficients is one term of the series; the Tauberian input is the size of the
whole series near the singularity.
\writeup (a) We have
\[
\begin{aligned}
\log P(e^{-t})&=-\sum_{n\ge1}\log(1-e^{-nt})=\sum_{n\ge1}\sum_{m\ge1}\frac{e^{-nmt}}m\\
&=\sum_{m\ge1}\frac1m\frac{e^{-mt}}{1-e^{-mt}}=\sum_{m\ge1}\frac1m\frac1{e^{mt}-1},
\end{aligned}
\]
all terms being
positive so the rearrangement is legitimate. Since $e^u-1\ge u$, each term is at most $\frac1{m^2t}$, and the sum is at most $\frac{\zeta(2)}t$. (b) Write
$g(u)=\frac u{e^u-1}$, so $0<g\le1$, $g$ is continuous and $g(0^+)=1$. Then $t\log P(e^{-t})=\sum_{m\ge1}\frac{g(mt)}{m^2}$, and each term tends to $\frac1{m^2}$ as
$t\to0$ while being dominated by $\frac1{m^2}$; by dominated convergence the sum tends to $\zeta(2)=\frac{\pi^2}6$. (c) All coefficients of $P$ are
non-negative, so $p(n)e^{-nt}\le P(e^{-t})$, giving $\log p(n)\le nt+\frac{\pi^2}{6t}$ for every $t>0$. The right side is minimal at
$t=\frac\pi{\sqrt{6n}}$, where it equals $2\sqrt{\frac{\pi^2n}6}=\pi\sqrt{\frac{2n}3}$. Note that Karamata's theorem is of no use here: $\log P$ is
regularly varying, but $P$ itself grows like $e^{c/t}$, faster than any power of $\frac1t$, so there is no $\rho$ to speak of. All the information used is
the single inequality $p(n)e^{-nt}\le P(e^{-t})$, which is why the bound is so cheap --- and why it cannot produce a lower bound of the same shape.
\end{solution}

\begin{solution}{ex:quicksort}
\idea Multiply the recurrence by $n$ to make the sum telescope, solve for the generating function, then let Karamata read off the sum of the $C_k$ and
monotonicity convert it into $C_n$.
\writeup (a) Since $H_n=\sum_{k\le n}\frac1k$ has $\sum_nH_nz^n=\frac1{1-z}\log\frac1{1-z}$, differentiating $\sum_nH_nz^{n+1}$ gives
$\sum_n(n+1)H_nz^n=\frac{L+z}{(1-z)^2}$ with $L=\log\frac1{1-z}$. One checks $C_n=2(n+1)H_n-4n$ satisfies the recurrence (both sides multiplied by $n$ give
$n(n-1)+2\sum_{k<n}C_k$), and $\sum_nnz^n=\frac z{(1-z)^2}$, so $C(z)=\frac{2(L+z)}{(1-z)^2}-\frac{4z}{(1-z)^2}=\frac{2(L-z)}{(1-z)^2}$. (b) From
$C_n=2(n+1)H_n-4n$, $C_{n+1}-C_n=2(n+2)H_{n+1}-2(n+1)H_n-4=2H_n+\frac{2(n+2)}{n+1}-4\ge0$ for $n\ge1$. (c) With $t\to0^+$ and $z=e^{-t}$ we have $1-z\sim t$ and
$L\sim\log\frac1t$, so $C(e^{-t})\sim\frac{2\log\frac1t}{t^2}=t^{-2}\ell(\frac1t)$ with $\ell(y)=2\log y$ slowly varying. By
Corollary~\ref{cor:karamataseries} with $\rho=2$, $S_n=\sum_{k\le n}C_k\sim\frac{2n^2\log n}{\Gamma(3)}=n^2\log n$. For $0<\theta<1$, monotonicity gives
$(n-\lfloor\theta n\rfloor)C_{\lfloor\theta n\rfloor}\le S_n-S_{\lfloor\theta n\rfloor}\le(n-\lfloor\theta n\rfloor)C_n$, and $S_n-S_{\theta n}\sim(1-\theta^2)n^2\log n$. The right
inequality gives $\liminf\frac{C_n}{n\log n}\ge\frac{1-\theta^2}{1-\theta}=1+\theta$; applying the left inequality with $n$ replaced by $\frac n\theta$ gives
$\limsup\frac{C_n}{n\log n}\le\frac{\theta^{-2}-1}{\theta^{-1}-1}=1+\frac1\theta$. Letting $\theta\to1^-$ gives $C_n\sim2n\log n$. (Quicksort on a million items makes about $2.5\cdot10^7$ comparisons on average (the exact $C_n=2(n+1)H_n-4n$), against
$\log_2(10^6!)\approx1.85\cdot10^7$ for the information-theoretic minimum: a third more. The leading terms $2n\log n$ and $n\log_2n$ have
ratio $2\ln2\approx1.39$, so the gap settles at $39\%$.)
\end{solution}

\begin{solution}{ex:wienerZ}
\idea On $\Z$ the local-to-global work was already done in Chapter~16: Wiener's lemma inverts $\hat a$ in one stroke, and then there is nothing to prove.
\writeup By Theorem~3.11 applied to $\frac1{\hat a}$ --- that is, by Wiener's lemma in $W$, transported to $\ell^1(\Z)$ by the identification of
Section~3.6 --- there is $b\in\ell^1(\Z)$ with $\hat b=\frac1{\hat a}$, so $a*b=\delta_0$. Let $y=a*x\in\ell^\infty$, which tends to $s$. Then
$x=\delta_0*x=b*a*x=b*y$ (convolution of $\ell^1$ with $\ell^\infty$ is associative, by absolute convergence), so $x_n=\sum_kb_ky_{n-k}$. Each term tends to
$b_ks$ and is dominated by $|b_k|\|y\|_\infty$, so by dominated convergence $x_n\to s\sum_kb_k=s\hat b(0)=\frac s{\hat a(0)}$. For
$a=(\frac23,\frac13)$ at lags $0,1$: $\hat a(\theta)=\frac{2+e^{i\theta}}3\ne0$ since $|e^{i\theta}|=1<2$, and $\hat a(0)=1$, so $x_n\to c$. For
$a=(\frac12,\frac12)$: $\hat a(\pi)=0$, and indeed $x_n=(-1)^n$ has $\frac{x_n+x_{n-1}}2=0$ for all $n$ while $x_n$ diverges. One vanishing frequency, one
invisible oscillation.
\end{solution}

\begin{solution}{ex:ikeharaline}
\idea A pole of $\zeta$ shifted up the line produces a term $x^{1+i}$ in the summatory function, which is as large as $x$ and oscillates.
\writeup (a) $\cos(\log n)=\frac{n^i+n^{-i}}2$, so $\sum_na_nn^{-s}=\zeta(s)+\frac12\sum_n(n^{-(s-i)}+n^{-(s+i)})$, and all three series converge absolutely for
$\operatorname{Re}s>1$. (b) $\sum_{n\le x}1=x+O(1)$. For the cosine, $u(t)=\cos(\log t)$ has $|u'(t)|=\frac{|\sin\log t|}t\le\frac1t$, so for each $n$,
$|u(n)-\int_n^{n+1}u|\le\frac1n$, whence $\sum_{n\le x}u(n)=\int_1^xu(t)\dd t+O(\log x)$. Now
$\int_1^xt^i\dd t=\frac{x^{1+i}-1}{1+i}$, and $\operatorname{Re}\frac{x^{1+i}}{1+i}=\operatorname{Re}\frac{x e^{i\log x}(1-i)}2=\frac x2(\cos\log x+\sin\log x)=\frac
x{\sqrt2}\cos(\log x-\frac\pi4)$. (c) $F(s)-\frac1{s-1}$ has poles at $s=1\pm i$, which lie \emph{on} the line $\operatorname{Re}s=1$, so it does not extend
continuously to the closed half-plane. In the proof of Theorem~\ref{thm:ikehara} the smoothing step \eqref{eq:ikesmooth} is still valid for $\varepsilon>0$, but
$H(1+\varepsilon+2\pi i\xi)$ blows up as $\varepsilon\to0$ near $\xi=\pm\frac1{2\pi}$, so no limit can be taken: the two frequencies $\pm\frac1{2\pi}$ are exactly the
oscillation $\cos(\log x)$ that survives in (b).
\end{solution}

\section*{Sources and further reading}
\addcontentsline{toc}{section}{Sources and further reading}

N.~H.~Abel, ``Untersuchungen \"uber die Reihe $1+\frac m1x+\dots$'', \emph{J.\ reine angew.\ Math.} 1 (1826), 311--339, contains Abel's continuity theorem
and the famous footnote on Cauchy's theorem; his letters of 1826, including the one to Holmboe complaining of the lack of rigour in the treatment of
series, are in \emph{\OE uvres compl\`etes}, vol.~2 (Christiania, 1881), 254--265. The sentence about divergent series being the devil's invention is
universally attributed to Abel but we have not found it in his published writings, so we have paraphrased his documented complaint instead.

A.~Tauber, ``Ein Satz aus der Theorie der unendlichen Reihen'', \emph{Monatsh.\ Math.\ Phys.} 8 (1897), 273--277. The biographical details (the Ph\"onix
insurance company, the deportation to Theresienstadt in June 1942) are from the MacTutor History of Mathematics archive, \emph{Alfred Tauber}, and from
C.~Binder's article in \emph{\"Osterreichische Mathematiker}. J.~E.~Littlewood, ``The converse of Abel's theorem on power series'', \emph{Proc.\ London
Math.\ Soc.} (2) 9 (1911), 434--448; his account of how long the proof took him is in \emph{A Mathematician's Miscellany} (Methuen, 1953). G.~H.~Hardy and
J.~E.~Littlewood, ``Tauberian theorems concerning power series and Dirichlet's series whose coefficients are positive'', \emph{Proc.\ London Math.\ Soc.}
(2) 13 (1914), 174--191, is where the adjective ``Tauberian'' was launched.

J.~Karamata, ``\"Uber die Hardy--Little\-wood\-schen Um\-keh\-rungen des Abel\-schen Stetig\-keits\-satzes'', \emph{Math.\ Z.} 32 (1930), 319--320: two pages. The theory
of regularly varying functions is developed in N.~H.~Bingham, C.~M.~Goldie and J.~L.~Teugels, \emph{Regular Variation} (Cambridge, 1987).

N.~Wiener, ``Tauberian theorems'', \emph{Ann.\ of Math.} 33 (1932), 1--100; the kernel-and-geometric-series proof of Lemma~\ref{lem:localinv} follows
W.~Rudin, \emph{Functional Analysis}, 2nd ed. (McGraw-Hill, 1991), Chapter~9, and H.~Reiter, \emph{Classical Harmonic Analysis and Locally Compact Groups}
(Oxford, 1968). S.~Ikehara, ``An extension of Landau's theorem in the analytic theory of numbers'', \emph{J.\ Math.\ Phys.\ (MIT)} 10 (1931), 1--12;
S.~Bochner's kernel proof is in his \emph{Vorlesungen \"uber Fouriersche Integrale} (Akademische Verlagsgesellschaft, 1932); the version we give follows
D.~J.~Newman's and J.~Korevaar's expositions.

E.~Landau, ``\"Uber die Einteilung der positiven ganzen Zahlen in vier Klassen nach der Mindestzahl der zu ihrer additiven Zusammensetzung erforderlichen
Quadrate'', \emph{Arch.\ Math.\ Phys.} (3) 13 (1908), 305--312, is the original proof of $B(x)\sim Kx/\sqrt{\log x}$; Ramanujan's announcement is in his
first letter to Hardy, 16 January 1913, reprinted in \emph{Collected Papers} (Cambridge, 1927), xxiii--xxx. H.~Delange, ``G\'en\'eralisation du
th\'eor\`eme de Ikehara'', \emph{Ann.\ Sci.\ \'Ec.\ Norm.\ Sup.} (3) 71 (1954), 213--242, handles singularities of the form $(s-1)^{-\alpha}$ and gives
Landau's form with an expansion. For the constant to many digits and its history, see S.~R.~Finch, \emph{Mathematical Constants} (Cambridge, 2003),
Section~2.3.

The two indispensable books are G.~H.~Hardy, \emph{Divergent Series} (Oxford, 1949), written by one of the two people who made the subject, and
J.~Korevaar, \emph{Tauberian Theory: A Century of Developments} (Springer, 2004), which is the modern survey. For the analysis of algorithms by these
methods, P.~Flajolet and R.~Sedgewick, \emph{Analytic Combinatorics} (Cambridge, 2009), Chapters~IV--VIII, and the quicksort recurrence in
D.~E.~Knuth, \emph{The Art of Computer Programming}, vol.~3, 2nd ed. (Addison-Wesley, 1998), Section~5.2.2.
